\begin{table*}[t]
\centering
\caption{Generation grid and task-family partition used in the reported experiment.}
\label{tab:dataset-spec}
\scriptsize
\begin{tabularx}{\textwidth}{p{0.22\textwidth}p{0.24\textwidth}X}
\toprule
Category & Parameter & Value \\
\midrule
Dataset size & Java classes & 1,232 classes \\
& Questions per model & 46,656: 11,664 unguarded and 34,992 condition-aware \\
& Question polarity & 23,328 reachable and 23,328 unreachable \\
\midrule
Graph generation & \path{class_size} & $\{100,250,500\}$ methods \\
& \path{depths} & $\{1,4,7,10,13,16,19,22,25\}$ \\
& Local tree structure & One structured execution tree spanning all generated methods; branching probability $0.65$ \\
& \path{useful_var_spawn_probability} & 0.2 \\
\midrule
Generated code & \path{n_comment} & $\{0,1\}$ comment lines per method \\
& \path{n_vars}, \path{n_loops}, \path{n_params} & $\{0,1,2\}$ \\
& \path{n_if} & $\{0,1,3,5\}$ \\
\midrule
Factorial design & Code configurations & $3\times2\times3\times3\times4\times3=648$ \\
& \path{n_questions} & Four reachable and four unreachable per distance and configuration \\
\midrule
Task semantics & Unguarded (\mbox{\texttt{n\_if}$=0$}) & No generated condition guards a call \\
& Condition-aware (\mbox{\texttt{n\_if}$>0$}) & Only executable calls, including calls whose guards evaluate to true, contribute to reachability \\
\bottomrule
\end{tabularx}
\end{table*}
